\sheettitle
  {Optional course extract}
  {The $L^p$ interfaces retained for the July first pass}

This extract is omitted from the main document by default because the preparatory notes already contain the proofs. It records only the statements that must remain available for later deployment.

\section*{Ambient object and probability hierarchy}

For $1\le p<\infty$,
\[
L^p(\mu)
=
\left\{f:\int|f|^p\,d\mu<\infty\right\}\big/\!\sim_\mu,
\qquad
f\sim_\mu g
\Longleftrightarrow
f=g\quad\mu\text{-almost everywhere}.
\]
If $\mu(E)<\infty$ and $1\le p<q\le\infty$, then
\[
\|f\|_p
\le
\mu(E)^{1/p-1/q}\|f\|_q.
\]
Therefore, on a probability space,
\[
L^\infty(\mathbb P)
\subseteq
L^2(\mathbb P)
\subseteq
L^1(\mathbb P).
\]

\section*{Product interfaces}

\[
Y\in L^1(\mu),\ Z\in L^\infty(\mu)
\Longrightarrow
YZ\in L^1(\mu),
\qquad
\|YZ\|_1\le\|Y\|_1\|Z\|_\infty,
\]
and
\[
U,V\in L^2(\mu)
\Longrightarrow
UV\in L^1(\mu),
\qquad
\|UV\|_1\le\|U\|_2\|V\|_2.
\]
The first estimate controls bounded tests; the second controls cross-moments.

\section*{Density generation and change of measure}

If $h\ge0$ is measurable, then $h\mu$ is a positive measure absolutely continuous with respect to $\mu$. If $f\in L^1(\mu)$, then
\[
\nu_f=f\mu
\]
is a finite signed measure and
\[
|\nu_f|=|f|\mu.
\]
Moreover,
\[
\nu_f=\nu_g
\Longrightarrow
f=g\quad\mu\text{-almost everywhere}.
\]
If $\nu=h\mu$, then
\[
\int\varphi\,d\nu
=
\int\varphi h\,d\mu
\]
for nonnegative measurable $\varphi$, and for signed $\varphi$ under absolute integrability.

\section*{Representation boundary}

Forward density generation requires no RN theorem. Conversely, in the standard positive RN regime,
\[
\rho\ll\mu,
\qquad
\rho,\mu\text{ are }\sigma\text{-finite}
\Longrightarrow
\rho=\frac{d\rho}{d\mu}\,\mu,
\]
with uniqueness $\mu$-almost everywhere. A finite signed numerator is represented by applying the positive theorem to its Jordan parts.
